%==============================================================================
\section{How every claim in this report was checked}\label{app:verify}
%==============================================================================

This appendix describes the verification applied to the results of this report, so
that a reader who wishes to check them can do so directly. Section~\ref{sec:ai}
refers to it.

\subsection{Principle}\label{ssec:verify-principle}

Three rules were followed.

\emph{Independence.} The verification program shares no code with the solvers, the
campaign harness, or any script written during the research. Agreement between them is
therefore corroboration, not a second run of one implementation.

\emph{No assumed results.} The program does not use the results it is checking. The
loading rule of Proposition~\ref{prop:ktns} is not assumed optimal; it is implemented
separately and compared against an exact dynamic program over magazine states, so the
proposition itself is tested. The optimum $\Zst$ is computed by a dynamic program over
subsets of jobs, which is cross-checked against brute-force enumeration of all
permutations on the smallest instances, and which uses none of the methods of
Section~\ref{sec:exact}.

\emph{Exhaustive where possible.} On small instances the program enumerates rather than
samples. Where a statement concerns a family too large to enumerate, the sample and its
generation parameters are stated, and the statement is labelled an Observation rather
than a Proposition.

\subsection{What is covered}\label{ssec:verify-covered}

\begin{itemize}[topsep=4pt,itemsep=2pt]
  \item \textbf{The loading oracle.} KTNS against the exact dynamic program, on
  $83{,}670$ instance-and-sequence pairs, including every permutation of every instance
  in a random bank. No disagreement. The same comparison was run against the
  implementation used by the solvers, with the same result.
  \item \textbf{The convention identity} of Proposition~\ref{prop:conv}, on random
  instances.
  \item \textbf{Both lower bounds} and Corollary~\ref{cor:lb}, on every instance in the
  random census.
  \item \textbf{Every named instance} in Sections~\ref{sec:prelim} and \ref{sec:gap}:
  the $k$-ring for $k=3,\dots,9$ in both readings; the sliding-window ring; the
  witness $W$ of Proposition~\ref{prop:refute}; the pair $I_0$, $I_1$ of
  Proposition~\ref{prop:noclutter}; the $\Kst=4$ instance; the smallest sub-optimal
  instance of Observation~\ref{obs:smallest}; the copies of
  Observation~\ref{obs:unbounded-true}; the star and non-matroid instances of
  Propositions~\ref{prop:chromK} and \ref{prop:nonmatroid}.
  \item \textbf{Every bound in Section~\ref{ssec:ratio} and \ref{ssec:law}}, checked
  for violations on a random census of $200$ instances: Theorem~\ref{thm:uncond},
  Corollaries~\ref{cor:zerogap}, \ref{cor:smallZ} and \ref{cor:z3},
  Lemma~\ref{lem:transcap}, and Propositions~\ref{prop:k3} and \ref{prop:genk}. No
  violation.
  \item \textbf{The closed form of Proposition~\ref{prop:hk3}}, on instances with
  $\Kst=3$ drawn at random until $45$ were obtained. No disagreement.
  \item \textbf{The covering relaxation} of Proposition~\ref{prop:oddring}, by solving
  the fractional cover as a linear program for $k=3,\dots,8$. This is what identified
  that the statement requires $k\ge4$.
  \item \textbf{The setup-cost collapse} of Theorem~\ref{thm:collapse}, by enumerating
  all feasible groupings and all orderings of an instance and minimising the augmented
  objective directly. This is also what identified that the threshold must be read in
  the walk value.
  \item \textbf{The PCF relaxation} of Proposition~\ref{prop:pcf}, by building the
  linear program explicitly over all configurations and solving it, with and without
  the counting rows.
\end{itemize}

\subsection{What the verification found}\label{ssec:verify-found}

The checking is described here rather than only its conclusion, because two of its
findings changed the content of this report.

The first is the distinction of Section~\ref{ssec:frames}. An earlier version of this
work reported the $k$-ring family as the canonical example of a positive gap for the
two-phase heuristic, and reported the gap as unbounded on copies of it. The
verification found that every $k$-ring has \emph{zero} gap for the method as specified,
and that the values previously reported were walk values. That is why
Section~\ref{ssec:frames} exists, why Table~\ref{tab:rings} has two gap columns, and
why Proposition~\ref{prop:unbounded} is now stated for the walk value with
Observation~\ref{obs:unbounded-true} supplying a different family for the other
reading.

The second is Proposition~\ref{prop:oddring}, which had been stated for all $k$ and
fails at $k=3$, where the whole job set forms a single feasible group.

Both were errors of statement rather than of proof: the arguments were sound for the
objects they were really about. That is the characteristic failure mode this appendix
exists to catch, and it is the reason the verification does not assume the results it
checks.

A third finding concerns the campaign analysis rather than the structural results, and
is recorded in Section~\ref{ssec:validation} and Appendix~\ref{app:repro}. An earlier
version of that analysis identified an instance by its family and file name. The Crama
collection reuses file names across four magazine capacities, so rows describing
different instances were merged, and the merge produced apparent disagreements between
methods. With the capacity included in the identifier the disagreements vanish: all
$41{,}673$ pairs agree. The same version also computed the coverage bound from the tool
count declared in the instance file, which is wrong on the ten instances that declare a
tool no job requires. Neither error was in a solver. Both were in the code that reads
the results, which is the part least likely to be checked and, on this evidence, the
part most worth checking.

\subsection{Running it}\label{ssec:verify-running}

The program is \texttt{verification/verify\_report\_independent.py} in the repository.
It requires Python~3 and SciPy, needs no commercial solver, and completes in roughly
twenty-five minutes. It prints one line per claim, naming the section of this report the
claim comes from, the value that section states, and the value the computation
produces, and exits with a non-zero status if any of them disagree.

The campaign numbers of Section~\ref{sec:comp} are produced by a separate analysis over
the raw per-instance result files, described in Appendix~\ref{app:repro}. No number in
Section~\ref{sec:comp} is carried over from an earlier campaign.
